
\subsubsection{Neural Functionals}

Zhou et al. (2023) introduced Neural Functional Transformers (NFTs), demonstrating that attention mechanisms can process neural network weights for downstream tasks. Their validation focused on homogeneous MNIST MLP collections. Zhou et al. (2024) developed Universal Neural Functionals with algorithms for constructing permutation-equivariant models for arbitrary architectures, though without large-scale empirical validation on heterogeneous model populations. Zhou et al. (2023) established the framework for permutation-equivariant neural functionals through NF-Layers with parameter sharing, demonstrating applications to generalization prediction and implicit neural representation tasks.

\subsubsection{Weight-Space Tokenization}

Schürholt et al. (2024) proposed SANE, which uses sequential weight chunking with transformer backbones for scalable weight-space learning. Their approach processes weight subsets sequentially without explicit architecture-specific handling. Eilertsen et al. (2020) released the NWS dataset containing 320K weight snapshots from 16K trained networks, establishing weight space as a research domain.

\subsubsection{Equivariant Weight Processing}

Kofinas et al. (2024) proposed representing neural networks as computational graphs and using graph neural networks to learn equivariant representations across diverse architectures, reporting improvements on implicit neural representation classification and generalization prediction tasks.

\subsubsection{Model Zoo Resources}

Falk et al. (2025) released the ViT Model Zoo containing 250 unique Vision Transformer models with diverse generating factors, extending model population methods to state-of-the-art architectures.

\subsubsection{Relationship to Current Work}

Prior work establishes that attention mechanisms are effective for weight-space learning and that architecture-aware processing captures structural information. The current work examines whether compositional design---architecture-specific tokenization followed by shared NFT processing---can enable cross-architecture learning on heterogeneous model zoos at proof-of-concept scale.
